\chapter{Estendendo a Toalha}
\label{cap:toalha}

{\bfseries estender a toalha -- o Outro contingente -- destinatários}

Três cartas a quem se destina -- se orienta -- este texto-rito;
    a banca de qualificação é o Outro contingente:
        Susana, Marta e Pedro são meus interlocutores neste ritual.

Primeiro veio Pedro, que me orientou inicialmente trazendo Latour ao meu olhar;
    depois encontrei Susana, com quem conflui desde as primeiras leituras, e me orientou com delicadeza e potência na realização deste ritual:
        e no começo, novamente, veio Marta, contra-Latour, cientista da práxis e da guerrilha

---

As cartas dirigir-se-ão aos três de modo a localizar a qualquer leitor quem é-são os leitores implícitos do texto;
    ou seja:
        qual máscara devem vestir para participar deste texto-rito.

---

{\bfseries Cartas contingentes:\\[+5pt]}
    \hspace*{2cm}. Pedro\\
    \hspace*{2cm}. Susana\\
    \hspace*{2cm}. Marta

Localizar também é localizar o interlocutor;
    grandes romances fazem isso, e as cartas são uma forma de informar localizando o interlocutor, dando-lhe um ponto de fuga de apoio que facilite a articulação do que se está sendo dito.

Os destinatários são os sacrificantes do rito, não as vítimas, nem os sacerdotes-feiticeiros;
    mas o rito é de passagem e intenciona transformar todos -- público e sacrificantes -- em feiticeiros.